\subsection{Core Language GADT}
\subsubsection{What is a GADT}
\evan{Bill, if you think we need this section I think you should be the one to write it because you have a much better understanding of these things than me}

\subsubsection{Core AST GADT}
\par We use GADTs to index our core expressions by their depth, 
so during translation \code{Surface a} becomes \code{Core n a} where $n$ is the depth. 
The depth of a core expression can be used to force well formedness at compile time. 
For example consider the following constructor, \codeln{(Union \ccons String -> Core n a -> Core n a -> Core n a)}
This tells us that the constructor called ``Union'' requires a string (the label) as well as two cores of the same depth.
With this specification, if the compiler sees the constructor being applied to two cores that are not the same depth, it will
throw a type error. We can force similar constraints for cons and concat as well.

    \codeln{(\cons) :: Ast n a -> Ast (n + 1) a -> Ast (n + 1) a}
    \codeln{(\concat) :: Ast (n + 1) a -> Ast (n + 1) a -> Ast (n + 1) a}

We that the cons constructor takes an argument on the left, and then an argument on the right that is greater in depth by 1.
The concat operator requires two cores of the same depth, 
but saying the depth is $(n + 1)$ forces the depth to be non-zero so that both arguments are lists than can be concattenated.